\chapter{Administración de Bases de Datos}

\section[Backups]{\emj{🧠}~Backups}

Un backup es inútil si no ha sido probado. La restauración es la parte importante, no el respaldo.
PostgreSQL ofrece dos niveles de backup: lógico (SQL/dump) y físico (copia de archivos del cluster).
La estrategia óptima combina backups físicos periódicos con WAL archiving para Point-In-Time Recovery.

\begin{teoriabox}{Tipos de backup en PostgreSQL}
\begin{itemize}
  \item \textbf{pg\_dump:} backup lógico de una sola BD. Genera SQL o formato binario comprimido. Portable entre versiones y SOs. No incluye roles globales.
  \item \textbf{pg\_dumpall:} backup lógico de todas las BDs + roles globales. No paralelizable.
  \item \textbf{pg\_basebackup:} copia física del cluster completo. Necesario para PITR. Requiere acceso al servidor. Más rápido que pg\_dumpall en BDs grandes.
  \item \textbf{WAL Archiving + PITR:} archivar los WAL segments permite restaurar a cualquier punto en el tiempo, no solo al momento del último backup.
\end{itemize}
Estrategia recomendada: \texttt{pg\_basebackup} semanal + WAL archiving continuo + \texttt{pg\_dump} diario para cada BD crítica.
\end{teoriabox}

\begin{warningbox}{Nunca confíes en un backup no probado}
Ejecuta restauraciones de prueba regularmente en un ambiente aislado.
Verifica que los datos sean consistentes y que las aplicaciones funcionen contra la BD restaurada.
Un backup corrupto o incompleto es igual a no tener backup.
\end{warningbox}

\section[Usuarios, roles y permisos]{\emj{🧠}~Usuarios, roles y permisos}

PostgreSQL usa \textbf{roles} como unidad principal de seguridad. Un rol con el atributo \texttt{LOGIN} funciona como usuario; sin él, es solo un grupo de permisos.
Los permisos son granulares: a nivel de BD, schema, tabla, columna y fila (Row Level Security).
El principio de mínimo privilegio dicta que cada rol tenga exactamente los permisos que necesita, ni más.

\begin{itemize}
  \item Roles con \texttt{LOGIN}: usuarios que se conectan
  \item Roles sin \texttt{LOGIN}: grupos de permisos que se asignan a otros roles
  \item \texttt{GRANT OPTION}: permite al receptor otorgar ese permiso a otros
  \item \texttt{pg\_hba.conf}: controla autenticación (qué IPs pueden conectarse y con qué método)
\end{itemize}

\section[Monitoreo básico]{\emj{🧠}~Monitoreo básico}

El monitoreo proactivo previene incidentes; el monitoreo reactivo solo los atiende después del daño.
Las vistas del sistema de PostgreSQL ofrecen visibilidad en tiempo real sin necesidad de herramientas externas.

\begin{itemize}
  \item \texttt{pg\_stat\_activity}: conexiones activas, estado, query en ejecución
  \item \texttt{pg\_stat\_user\_tables}: estadísticas por tabla (lecturas, escrituras, dead tuples)
  \item \texttt{pg\_locks}: locks activos y conflictos
  \item \texttt{pg\_stat\_statements}: historial de queries con tiempos de ejecución (extensión)
\end{itemize}

\begin{lstlisting}[caption={Backups con pg\_dump y restauración}]
-- Backup lógico (formato custom, comprimido)
pg_dump -U postgres -Fc -f backup_mibd_20260516.dump mibd

-- Backup lógico en SQL plano
pg_dump -U postgres -f backup_mibd.sql mibd

-- Restaurar desde formato custom
pg_restore -U postgres -d mibd_nueva backup_mibd_20260516.dump

-- Restaurar SQL plano
psql -U postgres -d mibd_nueva -f backup_mibd.sql

-- Backup de todas las bases + roles globales
pg_dumpall -U postgres -f cluster_completo.sql
\end{lstlisting}

\begin{lstlisting}[caption={Monitoreo --- conexiones, locks y queries lentas}]
-- Conexiones activas
SELECT pid, usename, state, query_start,
       LEFT(query, 100) AS query_corta
FROM   pg_stat_activity
WHERE  state != 'idle'
ORDER BY query_start;

-- Terminar conexión colgada
SELECT pg_terminate_backend(pid)
FROM   pg_stat_activity
WHERE  state = 'idle in transaction'
  AND  query_start < NOW() - INTERVAL '10 minutes';

-- Ver locks activos
SELECT pid, relation::regclass, mode, granted
FROM   pg_locks
WHERE  NOT granted;
\end{lstlisting}

\begin{lstlisting}[caption={pg\_stat\_statements --- queries más lentas}]
-- Activar extensión (requiere superuser, reinicio de PG)
CREATE EXTENSION IF NOT EXISTS pg_stat_statements;

-- Top 10 queries más lentas por tiempo promedio
SELECT LEFT(query, 120) AS query,
       calls,
       ROUND((total_exec_time / calls)::numeric, 2) AS avg_ms,
       ROUND(total_exec_time::numeric, 2) AS total_ms
FROM   pg_stat_statements
ORDER BY avg_ms DESC
LIMIT  10;

-- Resetear estadísticas
SELECT pg_stat_statements_reset();
\end{lstlisting}

\begin{entrevistabox}
\begin{enumerate}
  \item ¿Cuál es la diferencia entre pg\_dump y pg\_basebackup?
  \item ¿Cómo identificarías una query lenta en producción sin bajar el servidor?
  \item ¿Qué es el principio de mínimo privilegio aplicado a roles de base de datos?
\end{enumerate}
\end{entrevistabox}
